\subsection{Hilbert 空间中的基}\label{sec:bases-hilbert-spaces}

\subsubsection{基的存在性}\label{subsec:hilbert-basis-existence}

由线性代数知道, 每个 Euclidean 向量空间都有 Cartesian 基. 本小节将把这一结果推广到 Hilbert 空间. 需要以下定义.

\begin{Definition}
\label{def:unordered-hilbert-series-sum}
设 $H$ 为 Hilbert 空间, $I$ 为一个 (可能不可数的) 指标集. 设 $\{x_i\}_{i\in I}$ 为 $H$ 的一族元素. 若至多可数多个 $x_i$ 非零, 且对这些非零元素的任意枚举都有 $x=\sum_{j\in\mathbb{N}}x_{i(j)}$, 就说 $\sum_{i\in I}x_i=x$.
\end{Definition}

\begin{Remark}
\label{rem:unordered-series-countability}
注意, 虽然允许指标集可能不可数会给我们带来方便, 但要使这一收敛概念有意义, 就至多只能有可数多个非零元素. 为说明这一点, 注意一个实数级数若绝对收敛, 则对于自然数 $n$, 范数大于例如 $1/n$ 的项至多只有有限个. 因而它至多只有可数多个非零项. 上述收敛定义推广了实数级数的绝对收敛, 以下条件容易证明与该定义等价. 证明留给读者 (Problem~\ref{ex:hilbert-series-convergence-criteria}).
\end{Remark}

\begin{Lemma}
\label{lem:unordered-hilbert-series-convergence-criteria}
和 $x=\sum_{i\in I}x_i$ 存在的充分必要条件是以下任一条件成立:\begin{enumerate}
\item 对每个 $\epsilon>0$, 存在 $I$ 的有限子集 $J_\epsilon$, 使对任意满足 $J_\epsilon\subset J\subset I$ 的有限集 $J$, 有
\begin{equation}\label{eq:hilbert-series-finite-sum-approximation}
\left\|x-\sum_{i\in J}x_i\right\|<\epsilon.
\end{equation}

\item 对每个 $\epsilon>0$, 存在 $I$ 的有限子集 $J_\epsilon$, 使
\begin{equation}\label{eq:hilbert-series-finite-tail-cauchy-condition}
\left\|\sum_{i\in J}x_i\right\|<\epsilon
\end{equation}
对 $I$ 的任意满足 $J\cap J_\epsilon=\emptyset$ 的有限子集 $J$ 成立.\end{enumerate}
\end{Lemma}

以下关注正交元素的和. 有如下引理.

\begin{Lemma}
\label{lem:orthogonal-series-summability}
设 $\{x_i\}_{i\in I}$ 为 Hilbert 空间 $H$ 中一族两两正交的元素. 那么 $\sum_{i\in I}x_i$ 存在的充分必要条件是 $\sum_{i\in I}\|x_i\|^2<\infty$. 在此情形下, 还有
\begin{equation}\label{eq:orthogonal-series-pythagorean-norm}
\left\|\sum_{i\in I}x_i\right\|^2=\sum_{i\in I}\|x_i\|^2.
\end{equation}

\end{Lemma}

\begin{Proof}
对 $I$ 的任意有限子集 $J$, 利用 $\{x_i\}_{i\in I}$ 中元素两两正交这一事实, 得到
\begin{equation}\label{eq:finite-orthogonal-sum-norm-calculation}
\left\|\sum_{i\in J}x_i\right\|^2=\left(\sum_{i\in J}x_i,\sum_{l\in J}x_l\right)=\sum_{i\in J}(x_i,x_i)=\sum_{i\in J}\|x_i\|^2.
\end{equation}
其余由 Lemma~\ref{lem:unordered-hilbert-series-convergence-criteria} 得出.
\end{Proof}

\begin{Definition}
\label{def:orthonormal-family}
若对每个 $i\in I$ 都有 $\|x_i\|=1$, 则称 $H$ 中一族两两正交的元素 $\{x_i\}_{i\in I}$ \textbf{标准正交}.
\end{Definition}
